\subsection{Additive Noise Model}

We consider observations with additive Gaussian noise:
\begin{equation}
Y_i(t) = \mu_i(t) + \varepsilon_i(t)
\end{equation}
where $\mu_i(t) = \beta(t) S_i(t) \sum_j C_{ij} I_j(t)$ is the mean incidence for region $i$ at time $t$, and $\varepsilon_i(t) \sim \mathcal{N}(0, \sigma^2)$ are independent noise terms.

The observation vector is $\mathbf{Y} = [Y_1(0), Y_2(0), \ldots, Y_1(T-1), Y_2(T-1)]^T \in \mathbb{R}^{2T}$, with mean vector:
\begin{equation}
\boldsymbol{\mu}(\boldsymbol{\phi}) = [\mu_1(0), \mu_2(0), \ldots, \mu_1(T-1), \mu_2(T-1)]^T
\end{equation}

Under the additive noise model, the observations are independently distributed:
\begin{equation}
Y_i(t) \sim \mathcal{N}(\mu_i(t), \sigma^2)
\end{equation}

The likelihood function is:
\begin{equation}
f(\mathbf{Y}|\boldsymbol{\phi}) = \prod_{t=0}^{T-1} \prod_{i=1}^{2} \frac{1}{\sqrt{2\pi\sigma^2}} \exp\left(-\frac{(Y_i(t) - \mu_i(t))^2}{2\sigma^2}\right)
\end{equation}

Taking the logarithm:
\begin{align}
\ell(\boldsymbol{\phi}) &= \log f(\mathbf{Y}|\boldsymbol{\phi}) \\
&= \sum_{t=0}^{T-1} \sum_{i=1}^{2} \left[-\frac{1}{2}\log(2\pi\sigma^2) - \frac{(Y_i(t) - \mu_i(t))^2}{2\sigma^2}\right] \\
&= -T\log(2\pi\sigma^2) - \frac{1}{2\sigma^2}\sum_{t=0}^{T-1}\sum_{i=1}^{2} (Y_i(t) - \mu_i(t))^2
\end{align}

Taking the partial derivative with respect to parameter $\phi_k \in
\{\theta, S_1(0), I_1(0), S_2(0), I_2(0)\}$:
\begin{align*}
\frac{\partial \ell}{\partial \phi_k} &= -\frac{1}{2\sigma^2} \frac{\partial}{\partial \phi_k}\sum_{t=0}^{T-1}\sum_{i=1}^{2} (Y_i(t) - \mu_i(t))^2 \\
&= -\frac{1}{2\sigma^2} \sum_{t=0}^{T-1}\sum_{i=1}^{2} 2(Y_i(t) - \mu_i(t)) \cdot \left(-\frac{\partial \mu_i(t)}{\partial \phi_k}\right) \\
&= \frac{1}{\sigma^2} \sum_{t=0}^{T-1}\sum_{i=1}^{2} (Y_i(t) - \mu_i(t)) \frac{\partial \mu_i(t)}{\partial \phi_k}
\end{align*}

The Fisher information matrix is:
\begin{align*}
  [\mathbf{J}]_{kl} &= \mathbb{E}\left[\frac{\partial \ell}{\partial \phi_k}\frac{\partial \ell}{\partial \phi_l}\right] \\
  &= \mathbb{E}\left[\frac{1}{\sigma^4} \sum_{t,i} \sum_{s,j} (Y_i(t) - \mu_i(t))(Y_j(s) - \mu_j(s)) \frac{\partial \mu_i(t)}{\partial \phi_k}\frac{\partial \mu_j(s)}{\partial \phi_l}\right] \\
  &= \frac{1}{\sigma^4} \sum_{t,i} \sum_{s,j} \mathbb{E}[(Y_i(t) - \mu_i(t))(Y_j(s) - \mu_j(s))] \frac{\partial \mu_i(t)}{\partial \phi_k}\frac{\partial \mu_j(s)}{\partial \phi_l}
\end{align*}

Since the noise terms are independent:
\begin{equation*}
\mathbb{E}[(Y_i(t) - \mu_i(t))(Y_j(s) - \mu_j(s))] = \mathbb{E}[\varepsilon_i(t)\varepsilon_j(s)] = \begin{cases}
\sigma^2 & \text{if } (i,t) = (j,s) \\
0 & \text{if } (i,t) \neq (j,s)
\end{cases}
\end{equation*}

Therefore:
\begin{align*}
  [\mathbf{J}]_{kl} &= \frac{1}{\sigma^4}
  \sum_{t=0}^{T-1}\sum_{i=1}^{2} \sigma^2 \frac{\partial
    \mu_i(t)}{\partial \phi_k}\frac{\partial \mu_i(t)}{\partial
    \phi_l} \\
  %
  %
  %
  &= \frac{1}{\sigma^2} \sum_{t=0}^{T-1}\sum_{i=1}^{2} \frac{\partial
    \mu_i(t)}{\partial \phi_k}\frac{\partial \mu_i(t)}{\partial
    \phi_l}
\end{align*}
Lacking a closed form expression for the partials above, we find them
using the sensitivity equations.

If we define the Jacobian matrix:
\begin{equation}
  \mathbf{G} = \frac{\partial \boldsymbol{\mu}}{\partial
    \boldsymbol{\phi}} = \begin{bmatrix} \frac{\partial
      \mu_1(0)}{\partial \theta} & \frac{\partial \mu_1(0)}{\partial
      S_1(0)} & \cdots & \frac{\partial \mu_1(0)}{\partial I_2(0)} \\
    %
    %
    %
    \frac{\partial \mu_2(0)}{\partial \theta} & \frac{\partial
      \mu_2(0)}{\partial S_1(0)} & \cdots & \frac{\partial
      \mu_2(0)}{\partial I_2(0)} \\
    %
    %
    %   
    \vdots & \vdots & \ddots & \vdots \\
    %
    %
    %
    \frac{\partial \mu_2(T-1)}{\partial \theta} & \frac{\partial
      \mu_2(T-1)}{\partial S_1(0)} & \cdots & \frac{\partial
      \mu_2(T-1)}{\partial I_2(0)}
  \end{bmatrix} \in \mathbb{R}^{2T \times 5}
\end{equation}

Then the Fisher Information Matrix can be written compactly as:
\begin{equation}
\mathbf{J} = \frac{1}{\sigma^2} \mathbf{G}^T \mathbf{G}
\end{equation}

The CRLB can be computed numerically using Cholesky factorization as detailed above.
